\documentclass[12pt,a4paper]{article}
\usepackage[utf8]{inputenc}
\usepackage[T5]{fontenc}
\usepackage[vietnamese]{babel}
\usepackage[margin=1.5cm,top=1.4cm,bottom=1.4cm]{geometry}
\usepackage{amsmath,amssymb,array,tabularx,enumitem,fancyhdr,tikz,xcolor}
\usetikzlibrary{calc,angles,quotes,arrows.meta,babel,patterns}
\definecolor{ink}{RGB}{30,55,70}

\newcommand{\headerQuote}{%
  \ifcase\value{page}\relax
    \or Math gives us hope that every problem has a solution% Trang 1
    \or Math + You = Champion% Trang 2
    \or Where Frustration Becomes Inspiration% Trang 3
    \or Turning Confusion into Clarity% Trang 4
    \or Don't be an X, be a Y of mathematics% Trang 5
    \or Don't get mad. Learn to add% Trang 6
    \or Math gives us hope that every problem has a solution% Trang 7
    \or Math + You = Champion% Trang 8
    \or Where Frustration Becomes Inspiration% Trang 9
    \or Turning Confusion into Clarity% Trang 10
    \or Don't be an X, be a Y of mathematics% Trang 11
    \or Don't get mad. Learn to add% Trang 12
    \else Math gives us hope that every problem has a solution%
  \fi
}
\pagestyle{fancy}\fancyhf{}
\lhead{\small\textit{\headerQuote}}
\rhead{\small\textit{Người biên soạn: Dương Gia Bảo}}
\rfoot{\small Trang \thepage}
\renewcommand{\headrulewidth}{0.4pt}
\setlength{\headheight}{18pt}
\setlength{\parindent}{0pt}
\setlength{\parskip}{1.5pt}
\setlength{\tabcolsep}{4pt}
\renewcommand{\arraystretch}{1.12}
\newcommand{\parttitle}[1]{\par\vspace{2pt}{\large\bfseries\color{ink}#1}\par\vspace{1pt}}
\newcommand{\four}[4]{\par\noindent\begin{tabularx}{\linewidth}{@{}*{4}{>{\raggedright\arraybackslash}X}@{}}\textbf{A.} #1&\textbf{B.} #2&\textbf{C.} #3&\textbf{D.} #4\end{tabularx}}
\newcommand{\two}[4]{\par\noindent\begin{tabularx}{\linewidth}{@{}*{2}{>{\raggedright\arraybackslash}X}@{}}\textbf{A.} #1&\textbf{B.} #2\\[2pt]\textbf{C.} #3&\textbf{D.} #4\end{tabularx}}
\newcommand{\answerbox}{\hfill\textit{Đáp số:}\ \fbox{\rule{0pt}{12pt}\hspace{2.2cm}}}
\setlist[enumerate,1]{label=\textbf{Câu \arabic*.},leftmargin=1.4cm,labelsep=.15cm,itemsep=1.5pt,topsep=1.0pt,parsep=0pt}
\setlist[enumerate,2]{label=\alph*),leftmargin=.6cm,itemsep=1.0pt,topsep=1.0pt,parsep=0pt}

\begin{document}
\noindent
\begin{minipage}[c]{0.38\textwidth}
    \centering
    \textbf{LÀM TOÁN ĐỂ HẠNH PHÚC}\\[-2pt]
    \rule{3.2cm}{0.6pt}\\[4pt]
    \textbf{ĐỀ CHÍNH THỨC}
\end{minipage}\hfill
\begin{minipage}[c]{0.35\textwidth}
    \centering
    \textbf{ĐỀ ÔN THI GK1}\\
    \textbf{NĂM HỌC 2026 -- 2027}\\[2pt]
    \textbf{ĐỀ SỐ 07}\\
    \textbf{Môn: TOÁN -- Khối lớp 10}\\
    \textit{\small Thời gian làm bài: 90 phút\\(không kể thời gian phát đề)}
\end{minipage}

\vspace{0.2cm}
\noindent
\begin{minipage}[b]{0.78\textwidth}
    \textbf{Họ và tên:}\dotfill\textbf{SBD:}\makebox[3.5cm]{\dotfill}
\end{minipage}\hfill
\begin{minipage}[b]{0.20\textwidth}
    \hfill
    \setlength{\fboxrule}{0.8pt}%
    \fbox{\begin{minipage}{2.3cm}
        \centering
        \textbf{\small Mã đề thi}\\
        \textbf{0007}
    \end{minipage}}
\end{minipage}
\vspace{0.2cm}

\parttitle{PHẦN I. Trắc nghiệm nhiều lựa chọn}
\textit{Thí sinh trả lời từ câu 1 đến câu 12. Mỗi câu hỏi chỉ chọn một phương án đúng duy nhất.}
\begin{enumerate}
\item Trong các câu sau đây, câu nào là một mệnh đề toán học?
\four{Hà Nội là thủ đô của nước Việt Nam.}{Hãy giải phương trình này ngay đi!}{Tổng ba góc trong một tam giác bằng $180^\circ$.}{Số 10 có phải là số đẹp không?}

\item Trong các mệnh đề sau, mệnh đề nào là mệnh đề \textbf{sai}?
\four{$\forall x \in \mathbb{R},\; x^2 + 1 > 0$.}{$\exists x \in \mathbb{R},\; x^2 + 1 = 0$.}{Số $17$ là một số nguyên tố.}{Tam giác đều có ba cạnh bằng nhau.}

\item Cho tập hợp $X = \{a; b; c\}$. Số tập hợp con có đúng $2$ phần tử của tập hợp $X$ là
\four{$8$.}{$3$.}{$6$.}{$4$.}

\item Cho hai tập hợp $A = \{1; 2; 3; 4; 5\}$ và $B = \{2; 4; 6; 8\}$. Tập hợp hiệu $A \setminus B$ bằng
\four{$\{1; 3; 5\}$.}{$\{2; 4\}$.}{$\{6; 8\}$.}{$\{1; 2; 3; 4; 5; 6; 8\}$.}

\item Bất phương trình nào sau đây là bất phương trình bậc nhất hai ẩn $x, y$?
\four{$2x^2 + 3y \le 5$.}{$x + 2xy - 3y > 0$.}{$\dfrac{2}{x} + y < 1$.}{$3x - 4y + 7 \ge 0$.}

\item Hệ bất phương trình nào sau đây \textbf{không phải} là hệ bất phương trình bậc nhất hai ẩn?
\four{$\begin{cases} x - 2y \le 3 \\ 2x + y > 1 \end{cases}$.}{$\begin{cases} x + y \ge 0 \\ y - 3 \le 0 \end{cases}$.}{$\begin{cases} x + y > 2 \\ x - y^2 \le 1 \end{cases}$.}{$\begin{cases} 3x - y + 1 \ge 0 \\ x \ge 0,\; y \ge 0 \end{cases}$.}

\item Điểm nào sau đây thuộc miền nghiệm của bất phương trình $2x - 3y + 6 > 0$?
\four{$M(0; 3)$.}{$N(1; 2)$.}{$P(-4; 0)$.}{$Q(0; 2)$.}

\item Miền nghiệm của bất phương trình $x + 2y - 2 \le 0$ là nửa mặt phẳng (kể cả đường thẳng bờ $d: x + 2y - 2 = 0$) chứa điểm nào sau đây?
\four{$A(2; 2)$.}{$B(1; 3)$.}{$C(-1; 2)$.}{$O(0; 0)$.}

\item Giá trị của biểu thức lượng giác $\cos 120^\circ$ bằng
\four{$\dfrac{\sqrt{3}}{2}$.}{$-\dfrac{\sqrt{3}}{2}$.}{$\dfrac{1}{2}$.}{$-\dfrac{1}{2}$.}

\item Cho tam giác $ABC$ có các cạnh $BC = a, CA = b, AB = c$. Khẳng định nào sau đây là \textbf{đúng}?
\four{$a^2 = b^2 + c^2 - 2bc \cos A$.}{$a^2 = b^2 + c^2 + 2bc \cos A$.}{$a^2 = b^2 + c^2 - 2bc \sin A$.}{$a^2 = b^2 + c^2 - bc \cos A$.}

\item Cho góc $\alpha$ thỏa mãn $90^\circ < \alpha < 180^\circ$ và $\sin \alpha = \dfrac{3}{5}$. Giá trị của $\cos \alpha$ bằng
\four{$-\dfrac{4}{5}$.}{$\dfrac{4}{5}$.}{$-\dfrac{3}{4}$.}{$\dfrac{16}{25}$.}

\item Giá trị của biểu thức $P = \sin^2 35^\circ + \cos^2 35^\circ + 2\tan 45^\circ$ bằng
\four{$1$.}{$2$.}{$3$.}{$4$.}
\end{enumerate}

\parttitle{PHẦN II. Trắc nghiệm đúng/sai}
\textit{Thí sinh trả lời từ câu 13 đến câu 14. Trong mỗi ý a), b), c), d) ở mỗi câu, thí sinh chọn đúng hoặc sai.}
\begin{enumerate}[resume]
\item Cho hai tập hợp số $A = [-3; 2)$ và $B = (0; 5]$.
\begin{enumerate}
    \item Tập hợp $A$ chứa đúng $5$ số nguyên.
    \item Tập hợp $B$ là một nửa khoảng.
    \item Giao của hai tập hợp là khoảng $A \cap B = (0; 2)$.
    \item Hiệu của hai tập hợp là $A \setminus B = [-3; 0)$.
\end{enumerate}

\item Cho tam giác $ABC$ có độ dài hai cạnh $AB = 6$, $AC = 8$ và góc xen giữa $\widehat{A} = 60^\circ$.
\begin{enumerate}
    \item Diện tích tam giác $ABC$ được tính theo công thức $S = \dfrac{1}{2} AB \cdot AC \cdot \sin A$.
    \item Độ dài cạnh $BC$ bằng $2\sqrt{13}$.
    \item Diện tích tam giác $ABC$ bằng $12\sqrt{3}$.
    \item Bán kính đường tròn ngoại tiếp $R$ của tam giác $ABC$ bằng $\dfrac{\sqrt{39}}{3}$.
\end{enumerate}
\end{enumerate}

\parttitle{PHẦN III. Trắc nghiệm trả lời ngắn}
\textit{Thí sinh trả lời từ câu 15 đến câu 18. Viết kết quả vào ô quy định.}
\begin{enumerate}[resume]
\item Cho hai tập hợp $A = \{x \in \mathbb{N} \mid (2x^2 - 5x + 2)(x^2 - 4) = 0\}$ và $B = \{x \in \mathbb{Z} \mid -2 \le x < 3\}$. Tìm số phần tử của tập hợp $A \cup B$.
\answerbox

\item Cho hệ bất phương trình bậc nhất hai ẩn $\begin{cases} x \ge 0,\; y \ge 0 \\ x + y \le 4 \\ 2x + y \le 6 \end{cases}$.
Biết miền nghiệm của hệ là một miền đa giác. Tìm giá trị lớn nhất của biểu thức $F(x, y) = 3x + 2y$ trên miền nghiệm đó.
\answerbox

\item Cho tam giác $ABC$ có độ dài ba cạnh lần lượt là $a = 13$, $b = 14$ và $c = 15$. Tính bán kính đường tròn nội tiếp $r$ của tam giác $ABC$.
\answerbox

\item \begin{minipage}[t]{0.62\linewidth}
Một khinh khí cầu bay ở độ cao không đổi trên bầu trời. Từ hai vị trí quan sát $A$ và $B$ trên mặt đất bằng phẳng cách nhau $300\text{ m}$, hai người cùng nhìn thấy khinh khí cầu $C$ dưới các góc nâng lần lượt là $\widehat{CAB} = 45^\circ$ và $\widehat{CBA} = 60^\circ$ (hai điểm $A, B$ và hình chiếu $H$ của khinh khí cầu thẳng hàng, với $H$ nằm giữa $A$ và $B$). Hỏi độ cao của khinh khí cầu so với mặt đất bằng bao nhiêu mét? (\textit{Làm tròn kết quả đến hàng đơn vị}).
\vspace{4pt}

\answerbox
\end{minipage}\hfill
\begin{minipage}[t]{0.35\linewidth}
\centering
\begin{tikzpicture}[scale=0.78, every node/.style={scale=0.85}]
    \coordinate (A) at (0,0);
    \coordinate (B) at (4.5,0);
    % AH = CH / tan(45) = CH, BH = CH / tan(60) = CH / sqrt(3)
    % AH + BH = CH (1 + 1/sqrt(3)) = 4.5 => CH = 4.5 * sqrt(3)/(sqrt(3)+1) = 2.85
    % AH = 2.85, BH = 1.65
    \coordinate (H) at (2.85,0);
    \coordinate (C) at (2.85,2.85);

    % Ground
    \draw[thick] (-0.5,0) -- (5.0,0);
    \draw[thick, ink] (A) node[below] {$A$} -- (B) node[below] {$B$};
    \filldraw[black] (A) circle (1.5pt);
    \filldraw[black] (B) circle (1.5pt);
    \filldraw[black] (H) circle (1.5pt) node[below] {$H$};

    % Balloon lines
    \draw[dashed, blue!80!black, thick] (C) -- (H);
    \draw[thick, teal!80!black] (A) -- (C) -- (B);
    \filldraw[red!80!black] (C) circle (3pt) node[above] {$C$};

    % Right angle marker
    \draw (2.65,0) -- (2.65,0.2) -- (2.85,0.2);

    % Distance & Angles
    \node at (2.25,-0.4) {\small $300\text{ m}$};
    \node at (0.8,0.25) {\small $45^\circ$};
    \node at (3.8,0.3) {\small $60^\circ$};
    \node[right, blue!80!black] at (2.85, 1.5) {$h$};
\end{tikzpicture}
\end{minipage}
\end{enumerate}

\parttitle{PHẦN IV. Tự luận}
\textit{Thí sinh trình bày lời giải chi tiết từ câu 19 đến câu 22.}
\begin{enumerate}[resume]
\item \textit{(1,0 điểm)} Cho hai tập hợp $A = [m - 2; m + 3]$ và $B = (-1; 5]$ với $m$ là tham số thực. Tìm tất cả các giá trị của $m$ để tập hợp $A$ là tập con của tập hợp $B$.

\item \textit{(0,75 điểm)} Cho góc nhọn $\alpha$ thỏa mãn $0^\circ < \alpha < 90^\circ$ và $\tan \alpha = 2$. Tính giá trị của biểu thức:
\[
P = \dfrac{3\sin \alpha - 2\cos \alpha}{\sin \alpha + 4\cos \alpha}
\]

\item \begin{minipage}[t]{0.58\linewidth}
\textit{(0,75 điểm)} Để đo chiều cao của ngọn hải đăng $CH$ đứng thẳng trên một mỏm đá sát bờ biển, người ta chọn hai điểm quan sát $A$ và $B$ trên mặt đất phẳng sao cho ba điểm $A, B, H$ thẳng hàng theo thứ tự đó, với khoảng cách $AB = 50\text{ m}$. Từ $A$, góc nâng nhìn lên đỉnh ngọn hải đăng là $\widehat{CAH} = 32^\circ$. Từ $B$, góc nâng nhìn lên đỉnh là $\widehat{CBH} = 48^\circ$. Tính chiều cao $h = CH$ của ngọn hải đăng (\textit{làm tròn kết quả đến một chữ số thập phân}).
\end{minipage}\hfill
\begin{minipage}[t]{0.38\linewidth}
\centering
\begin{tikzpicture}[scale=0.8, every node/.style={scale=0.85}]
    \coordinate (A) at (0,0);
    \coordinate (B) at (2.5,0);
    \coordinate (H) at (4.7,0);
    \coordinate (C) at (4.7,2.8);

    \draw[thick] (-0.4,0) -- (5.2,0);
    \draw[thick, ink] (A) node[below] {$A$} -- (B) node[below] {$B$} -- (H) node[below] {$H$};
    \draw[dashed, blue!80!black, thick] (C) -- (H);
    \draw[thick, teal!80!black] (A) -- (C) -- (B);
    \filldraw[red!80!black] (C) circle (2.5pt) node[above] {$C$};
    \filldraw[black] (A) circle (1.5pt);
    \filldraw[black] (B) circle (1.5pt);
    \filldraw[black] (H) circle (1.5pt);

    \draw (4.45,0) -- (4.45,0.25) -- (4.7,0.25);
    \node at (1.25,-0.35) {\small $50\text{ m}$};
    \node at (0.8,0.25) {\small $32^\circ$};
    \node at (3.2,0.3) {\small $48^\circ$};
    \node[right, blue!80!black] at (4.7, 1.4) {$h$};
\end{tikzpicture}
\end{minipage}

\item \textit{(0,5 điểm)} Một xưởng sản xuất bánh trung thu truyền thống làm hai loại bánh: bánh nướng (loại I) và bánh dẻo (loại II).
\begin{itemize}[leftmargin=0.5cm]
    \item Để làm một chiếc bánh nướng cần $0{,}06\text{ kg}$ đường và $0{,}08\text{ kg}$ bột mì.
    \item Để làm một chiếc bánh dẻo cần $0{,}05\text{ kg}$ đường và $0{,}04\text{ kg}$ bột mì.
\end{itemize}
Xưởng có sẵn $15\text{ kg}$ đường và $16\text{ kg}$ bột mì dự trữ. Mỗi chiếc bánh nướng bán ra thu được lợi nhuận $40\,000$ đồng, mỗi chiếc bánh dẻo thu được lợi nhuận $30\,000$ đồng. Hỏi xưởng cần sản xuất bao nhiêu chiếc bánh mỗi loại để thu được tổng lợi nhuận lớn nhất, biết xưởng sản xuất chiếc nào tiêu thụ hết chiếc đó?
\end{enumerate}

\newpage
% =========================================================================
% BẢNG ĐÁP ÁN VÀ LỜI GIẢI CHI TIẾT
% =========================================================================
\parttitle{BẢNG ĐÁP ÁN TỔNG HỢP}

\vspace{0.2cm}
\textbf{1. Phần I: Trắc nghiệm nhiều lựa chọn}
\begin{center}
{\setlength{\tabcolsep}{2.5pt}\small
\begin{tabular}{|c|c|c|c|c|c|c|c|c|c|c|c|}
\hline
\textbf{Câu 1} & \textbf{Câu 2} & \textbf{Câu 3} & \textbf{Câu 4} & \textbf{Câu 5} & \textbf{Câu 6} & \textbf{Câu 7} & \textbf{Câu 8} & \textbf{Câu 9} & \textbf{Câu 10} & \textbf{Câu 11} & \textbf{Câu 12} \\
\hline
\textbf{C} & \textbf{B} & \textbf{B} & \textbf{A} & \textbf{D} & \textbf{C} & \textbf{B} & \textbf{D} & \textbf{D} & \textbf{A} & \textbf{A} & \textbf{C} \\
\hline
\end{tabular}}
\end{center}

\vspace{0.2cm}
\textbf{2. Phần II: Trắc nghiệm đúng/sai}
\begin{center}
\begin{tabular}{|c|c|c|c|c|}
\hline
\textbf{Câu hỏi} & \textbf{Ý a)} & \textbf{Ý b)} & \textbf{Ý c)} & \textbf{Ý d)} \\
\hline
\textbf{Câu 13} & \textbf{Đúng} & \textbf{Đúng} & \textbf{Đúng} & \textbf{Sai} \\
\hline
\textbf{Câu 14} & \textbf{Đúng} & \textbf{Đúng} & \textbf{Đúng} & \textbf{Sai} \\
\hline
\end{tabular}
\end{center}

\vspace{0.2cm}
\textbf{3. Phần III: Trắc nghiệm trả lời ngắn}
\begin{center}
\begin{tabular}{|c|c|c|c|}
\hline
\textbf{Câu 15} & \textbf{Câu 16} & \textbf{Câu 17} & \textbf{Câu 18} \\
\hline
\textbf{5} & \textbf{10} & \textbf{4} & \textbf{190} \\
\hline
\end{tabular}
\end{center}

\vspace{0.2cm}
\textbf{4. Phần IV: Tự luận}
\begin{center}
\begin{tabular}{|c|l|}
\hline
\textbf{Câu hỏi} & \textbf{Kết quả tóm tắt} \\
\hline
\textbf{Câu 19} & $m \in (1; 2]$ hay $1 < m \le 2$. \\
\hline
\textbf{Câu 20} & $P = \dfrac{2}{3}$. \\
\hline
\textbf{Câu 21} & Chiều cao ngọn hải đăng $h \approx 71{,}4\text{ m}$. \\
\hline
\textbf{Câu 22} & Sản xuất \textbf{125 chiếc bánh nướng} và \textbf{150 chiếc bánh dẻo}; lợi nhuận lớn nhất là $9\,500\,000$ đồng. \\
\hline
\end{tabular}
\end{center}

\vspace{0.4cm}
% =========================================================================
% HÌNH VẼ MIỀN NGHIỆM ĐỘC LẬP NGOÀI BẢNG CHO CÂU 16 VÀ CÂU 22
% =========================================================================
\parttitle{HÌNH VẼ BIỂU DIỄN MIỀN NGHIỆM (HỆ BẤT PHƯƠNG TRÌNH)}

\begin{minipage}[t]{0.48\textwidth}
\centering
\textbf{Hình 1: Miền nghiệm Câu 16}\\[6pt]
\begin{tikzpicture}[scale=0.8, >=Stealth]
    % Trục toạ độ
    \draw[line width=1.2pt, ->] (-0.8,0) -- (4.5,0) node[right, font=\normalsize\bfseries] {$x$};
    \draw[line width=1.2pt, ->] (0,-0.8) -- (0,5.0) node[above, font=\normalsize\bfseries] {$y$};
    \node[below left, font=\normalsize\bfseries] at (0,0) {$O$};

    % Đa giác nghiệm tô màu
    \fill[fill=ink!18] (0,0) -- (3,0) -- (2,2) -- (0,4) -- cycle;
    \draw[line width=1.8pt, blue!70!black] (0,0) -- (3,0) -- (2,2) -- (0,4) -- cycle;

    % Các đường thẳng biên
    % d1: x + y = 4 (cắt trục tại (4;0), (0;4))
    \draw[line width=1.1pt, dashed, red!80!black] (-0.3, 4.3) -- (4.3, -0.3) node[below right, font=\small] {$d_1: x+y=4$};
    % d2: 2x + y = 6 (cắt trục tại (3;0), (0;6))
    \draw[line width=1.1pt, dashed, teal!80!black] (-0.2, 6.4) -- (3.3, -0.6) node[above right, font=\small] {$d_2: 2x+y=6$};

    % Tọa độ các đỉnh
    \filldraw[black] (0,0) circle (3pt);
    \filldraw[black] (3,0) circle (3.5pt) node[below, font=\normalsize\bfseries] {$A(3;0)$};
    \filldraw[black] (2,2) circle (3.5pt) node[above right, font=\normalsize\bfseries] {$B(2;2)$};
    \filldraw[black] (0,4) circle (3.5pt) node[left, font=\normalsize\bfseries] {$C(0;4)$};

    % Nét đứt từ B xuống trục
    \draw[dotted, thick] (2,0) node[below, font=\small] {$2$} -- (2,2) -- (0,2) node[left, font=\small] {$2$};
\end{tikzpicture}
\end{minipage}\hfill
\begin{minipage}[t]{0.48\textwidth}
\centering
\textbf{Hình 2: Miền nghiệm Câu 22 (Quy hoạch tuyến tính)}\\[6pt]
\begin{tikzpicture}[x=0.012cm, y=0.0095cm, >=Stealth]
    % Trục toạ độ
    \draw[line width=1.2pt, ->] (-25,0) -- (240,0) node[right, font=\normalsize\bfseries] {$x$};
    \draw[line width=1.2pt, ->] (0,-25) -- (0,340) node[above, font=\normalsize\bfseries] {$y$};
    \node[below left, font=\normalsize\bfseries] at (0,0) {$O$};

    % Đa giác nghiệm O(0;0), A(200;0), B(125;150), C(0;300)
    \fill[fill=ink!18] (0,0) -- (200,0) -- (125,150) -- (0,300) -- cycle;
    \draw[line width=1.8pt, blue!70!black] (0,0) -- (200,0) -- (125,150) -- (0,300) -- cycle;

    % Đường biên d1: 6x + 5y = 1500 (cắt x tại 250, y tại 300)
    \draw[line width=1.1pt, dashed, red!80!black] (-10, 312) -- (235, 18) node[above right, font=\small] {$d_1: 6x+5y=1500$};
    % Đường biên d2: 2x + y = 400 (cắt x tại 200, y tại 400)
    \draw[line width=1.1pt, dashed, teal!80!black] (-10, 420) -- (210, -20) node[below right, font=\small] {$d_2: 2x+y=400$};

    % Các đỉnh
    \filldraw[black] (0,0) circle (3pt);
    \filldraw[black] (200,0) circle (3.5pt) node[below, font=\normalsize\bfseries] {$A(200;0)$};
    \filldraw[black] (125,150) circle (3.5pt) node[above right, font=\normalsize\bfseries] {$B(125;150)$};
    \filldraw[black] (0,300) circle (3.5pt) node[left, font=\normalsize\bfseries] {$C(0;300)$};

    % Nét đứt dóng tọa độ đỉnh B
    \draw[dotted, thick] (125,0) node[below, font=\small] {$125$} -- (125,150) -- (0,150) node[left, font=\small] {$150$};
\end{tikzpicture}
\end{minipage}

\newpage
% =========================================================================
% HƯỚNG DẪN GIẢI CHI TIẾT
% =========================================================================
\parttitle{HƯỚNG DẪN GIẢI CHI TIẾT}

\textbf{PHẦN I. Trắc nghiệm nhiều lựa chọn}
\begin{itemize}[leftmargin=0.6cm, itemsep=3pt]
    \item \textbf{Câu 1:} Mệnh đề toán học là khẳng định mang tính đúng hoặc sai về một đối tượng toán học. Câu ``Tổng ba góc trong một tam giác bằng $180^\circ$'' là mệnh đề toán học (đúng). Chọn \textbf{C}.
    \item \textbf{Câu 2:} Mệnh đề $\exists x \in \mathbb{R},\; x^2 + 1 = 0$ là sai vì $x^2 + 1 \ge 1 > 0,\; \forall x \in \mathbb{R}$. Chọn \textbf{B}.
    \item \textbf{Câu 3:} Số tập con có đúng 2 phần tử của tập hợp $X = \{a; b; c\}$ gồm $\{a; b\}, \{a; c\}, \{b; c\}$, có 3 tập hợp con. Chọn \textbf{B}.
    \item \textbf{Câu 4:} $A \setminus B = \{x \in A \mid x \notin B\} = \{1; 3; 5\}$. Chọn \textbf{A}.
    \item \textbf{Câu 5:} $3x - 4y + 7 \ge 0$ là bất phương trình bậc nhất hai ẩn có dạng $ax + by + c \ge 0$. Chọn \textbf{D}.
    \item \textbf{Câu 6:} Hệ ở phương án C chứa $y^2$ nên không phải là hệ bất phương trình bậc nhất hai ẩn. Chọn \textbf{C}.
    \item \textbf{Câu 7:} Thay tọa độ $N(1; 2)$ vào: $2(1) - 3(2) + 6 = 2 > 0$ (thỏa mãn). Chọn \textbf{B}.
    \item \textbf{Câu 8:} Thay $O(0; 0)$ vào vế trái: $0 + 2(0) - 2 = -2 \le 0$ (thỏa mãn). Vậy miền nghiệm chứa điểm $O(0; 0)$. Chọn \textbf{D}.
    \item \textbf{Câu 9:} $\cos 120^\circ = -\cos(180^\circ - 120^\circ) = -\cos 60^\circ = -\dfrac{1}{2}$. Chọn \textbf{D}.
    \item \textbf{Câu 10:} Theo định lý côsin trong tam giác: $a^2 = b^2 + c^2 - 2bc \cos A$. Chọn \textbf{A}.
    \item \textbf{Câu 11:} Vì $90^\circ < \alpha < 180^\circ$ nên $\cos \alpha < 0$. Do đó:
    $\cos \alpha = -\sqrt{1 - \sin^2 \alpha} = -\sqrt{1 - \dfrac{9}{25}} = -\dfrac{4}{5}$. Chọn \textbf{A}.
    \item \textbf{Câu 12:} $P = (\sin^2 35^\circ + \cos^2 35^\circ) + 2\tan 45^\circ = 1 + 2(1) = 3$. Chọn \textbf{C}.
\end{itemize}

\vspace{0.2cm}
\textbf{PHẦN II. Trắc nghiệm đúng/sai}
\begin{itemize}[leftmargin=0.6cm, itemsep=3pt]
    \item \textbf{Câu 13:}
    \begin{itemize}
        \item a) Các số nguyên thuộc $A = [-3; 2)$ là $\{-3; -2; -1; 0; 1\}$, có đúng $5$ số nguyên $\Rightarrow$ \textbf{Đúng}.
        \item b) Tập hợp $B = (0; 5]$ là một nửa khoảng $\Rightarrow$ \textbf{Đúng}.
        \item c) Giao của hai tập hợp là $A \cap B = [-3; 2) \cap (0; 5] = (0; 2)$ $\Rightarrow$ \textbf{Đúng}.
        \item d) Hiệu $A \setminus B = [-3; 2) \setminus (0; 5] = [-3; 0]$ (vì $0 \in A$ và $0 \notin B$ nên $0 \in A \setminus B$). Mệnh đề ghi $[-3; 0)$ là $\Rightarrow$ \textbf{Sai}.
    \end{itemize}

    \item \textbf{Câu 14:}
    \begin{itemize}
        \item a) Công thức diện tích tam giác là $S = \dfrac{1}{2} AB \cdot AC \cdot \sin A$ $\Rightarrow$ \textbf{Đúng}.
        \item b) Áp dụng định lý côsin: $BC^2 = AB^2 + AC^2 - 2 AB \cdot AC \cdot \cos 60^\circ = 6^2 + 8^2 - 2 \cdot 6 \cdot 8 \cdot \dfrac{1}{2} = 52 \implies BC = \sqrt{52} = 2\sqrt{13}$ $\Rightarrow$ \textbf{Đúng}.
        \item c) $S = \dfrac{1}{2} \cdot 6 \cdot 8 \cdot \sin 60^\circ = 24 \cdot \dfrac{\sqrt{3}}{2} = 12\sqrt{3}$ $\Rightarrow$ \textbf{Đúng}.
        \item d) Theo định lý sin: $2R = \dfrac{BC}{\sin A} = \dfrac{2\sqrt{13}}{\sqrt{3}/2} = \dfrac{4\sqrt{13}}{\sqrt{3}} \implies R = \dfrac{2\sqrt{39}}{3}$. Mệnh đề ghi $R = \dfrac{\sqrt{39}}{3}$ là $\Rightarrow$ \textbf{Sai}.
    \end{itemize}
\end{itemize}

\vspace{0.2cm}
\textbf{PHẦN III. Trắc nghiệm trả lời ngắn}
\begin{itemize}[leftmargin=0.6cm, itemsep=3pt]
    \item \textbf{Câu 15:} Ta có:
    \[
    (2x^2 - 5x + 2)(x^2 - 4) = 0 \iff \left[\begin{array}{l} 2x^2 - 5x + 2 = 0 \iff x = 2 \text{ hoặc } x = \dfrac{1}{2} \\ x^2 - 4 = 0 \iff x = 2 \text{ hoặc } x = -2 \end{array}\right.
    \]
    Vì $x \in \mathbb{N}$ nên $A = \{2\}$.
    Tập hợp $B = \{x \in \mathbb{Z} \mid -2 \le x < 3\} = \{-2; -1; 0; 1; 2\}$.
    Do đó $A \cup B = \{-2; -1; 0; 1; 2\}$, có tất cả $5$ phần tử.
    \item \textbf{Câu 16:} Miền nghiệm của hệ là miền tứ giác $OABC$ với các đỉnh:
    $O(0; 0)$, $A(3; 0)$, $B(2; 2)$, $C(0; 4)$.
    Tính giá trị $F(x, y) = 3x + 2y$ tại các đỉnh:
    \begin{itemize}
        \item $F(O) = 0$;
        \item $F(A) = 3(3) + 2(0) = 9$;
        \item $F(B) = 3(2) + 2(2) = 10$;
        \item $F(C) = 3(0) + 2(4) = 8$.
    \end{itemize}
    Vậy giá trị lớn nhất là $10$ (đạt tại đỉnh $B(2; 2)$).
    \item \textbf{Câu 17:} Nửa chu vi tam giác $p = \dfrac{13 + 14 + 15}{2} = 21$.
    Diện tích tam giác theo công thức Heron:
    \[
    S = \sqrt{21(21 - 13)(21 - 14)(21 - 15)} = \sqrt{21 \cdot 8 \cdot 7 \cdot 6} = \sqrt{7056} = 84.
    \]
    Bán kính đường tròn nội tiếp: $r = \dfrac{S}{p} = \dfrac{84}{21} = 4$.
    \item \textbf{Câu 18:} Gọi $H$ là hình chiếu vuông góc của khinh khí cầu $C$ lên mặt đất. Trong hai tam giác vuông $\triangle CAH$ và $\triangle CBH$:
    \[
    AH = \dfrac{CH}{\tan 45^\circ} = h, \quad BH = \dfrac{CH}{\tan 60^\circ} = \dfrac{h}{\sqrt{3}}.
    \]
    Vì $H$ nằm giữa $A$ và $B$ nên:
    \[
    AB = AH + BH = h\left(1 + \dfrac{1}{\sqrt{3}}\right) = 300 \implies h = \dfrac{300}{1 + \dfrac{1}{\sqrt{3}}} = \dfrac{300\sqrt{3}}{\sqrt{3}+1} = 150(3 - \sqrt{3}) \approx 190{,}19\text{ m}.
    \]
    Làm tròn đến hàng đơn vị, độ cao là $190\text{ m}$.
\end{itemize}

\vspace{0.2cm}
\textbf{PHẦN IV. Tự luận}
\begin{itemize}[leftmargin=0.6cm, itemsep=4pt]
    \item \textbf{Câu 19:}
    Để tập hợp $A = [m - 2; m + 3]$ là tập con của tập hợp $B = (-1; 5]$, ta cần và đủ:
    \[
    \begin{cases}
    m - 2 > -1 \\
    m + 3 \le 5
    \end{cases}
    \iff
    \begin{cases}
    m > 1 \\
    m \le 2
    \end{cases}
    \iff 1 < m \le 2.
    \]
    Vậy tất cả các giá trị của $m$ thỏa mãn là $m \in (1; 2]$.

    \item \textbf{Câu 20:}
    Vì $0^\circ < \alpha < 90^\circ$ nên $\cos \alpha \ne 0$. Chia cả tử và mẫu của biểu thức $P$ cho $\cos \alpha$:
    \[
    P = \dfrac{\dfrac{3\sin \alpha - 2\cos \alpha}{\cos \alpha}}{\dfrac{\sin \alpha + 4\cos \alpha}{\cos \alpha}} = \dfrac{3\tan \alpha - 2}{\tan \alpha + 4}.
    \]
    Thay $\tan \alpha = 2$ vào biểu thức:
    \[
    P = \dfrac{3(2) - 2}{2 + 4} = \dfrac{4}{6} = \dfrac{2}{3}.
    \]

    \item \textbf{Câu 21:}
    Xét tam giác $\triangle ABC$, ta có $\widehat{CAB} = 32^\circ$ và $\widehat{CBH} = 48^\circ$ là góc ngoài tại đỉnh $B$ của $\triangle ABC$.
    Do đó:
    \[
    \widehat{ACB} = \widehat{CBH} - \widehat{CAB} = 48^\circ - 32^\circ = 16^\circ.
    \]
    Áp dụng định lý sin trong tam giác $\triangle ABC$:
    \[
    \dfrac{BC}{\sin \widehat{CAB}} = \dfrac{AB}{\sin \widehat{ACB}} \implies BC = \dfrac{AB \cdot \sin 32^\circ}{\sin 16^\circ} = \dfrac{50 \cdot \sin 32^\circ}{\sin 16^\circ} \approx 96{,}126\text{ m}.
    \]
    Trong tam giác vuông $\triangle CBH$ vuông tại $H$:
    \[
    h = CH = BC \cdot \sin \widehat{CBH} = 96{,}126 \cdot \sin 48^\circ \approx 71{,}439\text{ m}.
    \]
    Làm tròn đến một chữ số thập phân, chiều cao ngọn hải đăng là $71{,}4\text{ m}$.

    \item \textbf{Câu 22:}
    Gọi $x, y$ lần lượt là số chiếc bánh nướng và bánh dẻo mà xưởng sản xuất ($x, y \in \mathbb{N}$).
    \begin{itemize}
        \item Lượng đường sử dụng: $0{,}06x + 0{,}05y \le 15 \iff 6x + 5y \le 1500$.
        \item Lượng bột mì sử dụng: $0{,}08x + 0{,}04y \le 16 \iff 2x + y \le 400$.
    \end{itemize}
    Ta có hệ bất phương trình ràng buộc:
    \[
    \begin{cases}
    x \ge 0,\; y \ge 0 \\
    6x + 5y \le 1500 \\
    2x + y \le 400
    \end{cases}
    \]
    Miền nghiệm của hệ là miền tứ giác $OABC$ với các đỉnh:
    \begin{itemize}
        \item $O(0; 0)$;
        \item $A(200; 0)$ (giao điểm của $2x + y = 400$ với trục $Ox$);
        \item $B(125; 150)$ (nghiệm của hệ $\begin{cases} 6x + 5y = 1500 \\ 2x + y = 400 \end{cases}$);
        \item $C(0; 300)$ (giao điểm của $6x + 5y = 1500$ với trục $Oy$).
    \end{itemize}
    Tổng lợi nhuận thu được là $F(x, y) = 40x + 30y$ (nghìn đồng). Tính giá trị của $F$ tại các đỉnh:
    \begin{itemize}
        \item $F(O) = 0$;
        \item $F(A) = 40(200) + 30(0) = 8000$ (nghìn đồng) = $8$ triệu đồng;
        \item $F(B) = 40(125) + 30(150) = 5000 + 4500 = 9500$ (nghìn đồng) = $9{,}5$ triệu đồng;
        \item $F(C) = 40(0) + 30(300) = 9000$ (nghìn đồng) = $9$ triệu đồng.
    \end{itemize}
    Giá trị lớn nhất của $F(x, y)$ đạt được tại đỉnh $B(125; 150)$.\\
    Vậy xưởng cần sản xuất \textbf{125 chiếc bánh nướng} và \textbf{150 chiếc bánh dẻo} để thu được lợi nhuận lớn nhất là \textbf{9.500.000 đồng}.
\end{itemize}

\end{document}
